\documentclass{article}
\usepackage{graphicx} % Required for inserting images
\usepackage{fullpage}
\usepackage{mpemath}
\usepackage{dsfont}
\usepackage{amsmath}
\usepackage{algorithm,algorithmicx,algpseudocode}

\newcommand{\red}[1]{\textcolor{red}{#1}}


\title{Fast CVaR}
\author{Palash Agrawal, Maeve Burwell, Gersi Doko, Marek Petrik}
\date{}

\begin{document}

\maketitle

\section{Introduction}

Palash

https://arxiv.org/pdf/2310.07224

\section{Preliminaries}

Before describing CVaR and its solution we define the basic notation we use in the paper. 
We use calligraphic letters to denote sets and a tilde to denote random variables. 
We also adopt the standard convention that $\mathcal{A}^{\mathcal{B}}$ represents the set of all functions from a set $\mathcal{B}$ to a set $\mathcal{A}$ and 
treat vectors as a function from indexes to real numbers. $\mathbf{1}$ denotes the vector of ones and $\mathds{1}(bool)$ denotes the indicator function. 
Finally $\Delta^n$ denotes the standard $n$-simplex while $\Delta_k^n$ denotes the k-capped simplex \red{citation needed}.

\section{Formal Model}

\begin{equation}
  \label{eq:primal_cvar}
  CVaR_\alpha(\tilde{x}) = 
  \begin{mprog}
    \minimize{\xi \in \Real^n} \xi\tr \tilde{x} 
    \stc \xi \leq \frac{1}{1-\alpha} p
    \cs \xi \in \Delta^n
  \end{mprog}
  =
  \begin{mprog}
    \minimize{\xi \in \Real^n} \xi\tr \tilde{x} 
    \stc \xi \leq \frac{1}{1-\alpha} p
    \cs \xi\tr \mathbf{1} = 1
    \cs \xi \geq 0
  \end{mprog}
\end{equation}

\begin{equation}
  \label{eq:k_capped_simplex}
  \psi(x, k) = 
  \begin{mprog}
    \minimize{\xi \in \Real^n} \lvert \lvert \xi - x \rvert \rvert_2^2
    \stc \xi \in \Delta_k^n
  \end{mprog}
  =
  \begin{mprog}
    \minimize{\xi \in \Real^n} \lvert \lvert \xi - x \rvert \rvert_2^2
    \stc \xi\tr \mathbf{1} \leq k
    \cs \xi \geq 0   
  \end{mprog}
\end{equation}

\section{Algorithm}

\subsection{Piece-wise Linear Minimization}

\red{Gersi doesn't know what to write here.}

\subsection{Quick Quantile}

\begin{algorithm}
  \floatname{algorithm}{Quick Quantile}
  \algrenewcommand\algorithmicrequire{\textbf{Input: Random variable $\tilde{X}$ with domain $x$ and probability mass function $p$ and $\alpha \in [0,1]$}}
  \algrenewcommand\algorithmicensure{\textbf{Output: $\alpha$-quantile of $\tilde{X}$}}
  \caption{}\label{alg:quick_quantile}
  \begin{algorithmic}[1]
    \Require 
    \Ensure 
    \If{$n = 1$}
    \State \Return $x$
    \EndIf
    \State $pivot \gets rand(1:n)$
    \State $tail \gets \sum_{i=0}^n \mathds{1}(x[i] \leq x[pivot]) p[i]$
    \If{$\alpha \leq tail$}
    \State \Return $qq(vals[\mathds{1}(x[i] \leq x[pivot])], p[\mathds{1}(x[i] \leq x[pivot])], \alpha)$
    \Else
    \State \Return $qq(vals[\mathds{1}(x[i] > x[pivot])], p[\mathds{1}(x[i] > x[pivot])], \alpha - tail)$
    \EndIf
  \end{algorithmic}
\end{algorithm}

\subsection{Quick CVaR}

\begin{algorithm}
  \floatname{algorithm}{Quick CVaR}
  \algrenewcommand\algorithmicrequire{\textbf{Input: Random variable $\tilde{X}$ with probability mass function p and $\alpha \in [0,1]$}}
  \algrenewcommand\algorithmicensure{\textbf{Output: $CVaR_\alpha(\tilde{X})$}}
  \caption{}\label{alg:quick_cvar}
  \begin{algorithmic}[1]
    \Require 
    \Ensure 
    \State $q_\alpha = \text{quick\_quantile}(\tilde{X}, p, \alpha)$ 
    \State \Return $\sum_{i=0}^n \mathds{1}( x[i] \leq q_\alpha ) x[i] p[i]$
  \end{algorithmic}
\end{algorithm}


\end{document}
